\documentclass[12pt]{article}
\usepackage[utf8]{inputenc}
\usepackage[T1]{fontenc}
\usepackage{hyperref}
\title{Wewaw Language Resources for K-12 Curriculum Development}
\author{Compiled by Claude Sonnet 5 (claude-sonnet-5) for Andr\'as Kornai}
\date{}
\begin{document}
\maketitle

\section{Introduction}
Wewaw (ISO 639-3 wea) is described as an extinct Sino-Tibetan language of Myanmar. The search found nothing on a language by that name, and no resources of any kind. The only near match was ``Wewew'', listed in some sources as a divergent dialect of Sgaw Karen. Whether Wewew is the same as Wewaw, or is extinct, could not be established from these results, so no link between them is claimed.

\section{Primary Educational Resources}
None found.

\section{Linguistic and Reference Materials}
None found. Checking Glottolog (glottocode wewa1238) and the ISO 639-3 registry would be the next step.

\section{Religious and Community Resources}
None found.

\section{Digital and Audio Resources}
None found.

\section{Curriculum Development Context}
No curriculum is feasible from the evidence found. This dossier records only that the search came up empty.
\end{document}